% =====================================================================
%  Chapter 2 - Experimental Setup and Methodology
% =====================================================================

\chapter{Experimental Setup and Methodology}
\label{ch:setup}

This chapter describes the experimental framework on which the research work is built.
It briefly covers the hardware platforms, the sensors, the software stack and the data
used to design and evaluate the algorithms presented in the following chapters.

Section~\ref{sec:av24} and Section~\ref{sec:eav24} introduce the two autonomous
racecars: the Dallara AV-24, which races in the Indy Autonomous Challenge (IAC), and
the Dallara EAV-24, which races in the Abu Dhabi Autonomous Racing League (A2RL).
Section~\ref{sec:sensors} presents the working principle, the advantages and the
limitations of the exteroceptive sensors (LiDARs, RaDARs and cameras) used for opponent
tracking, together with the sensor configuration of the two vehicles. Section~\ref{sec:sw_arch} gives an overview of the
main modules of the software architecture. Section~\ref{sec:datasets} describes the two
datasets used for the analysis and Section~\ref{sec:gt} explains how the opponent ground
truth is reconstructed.

% =====================================================================
\section{Dallara AV-24}
\label{sec:av24}

The Dallara AV-24 (Fig.~\ref{fig:av24}) is the official vehicle of the IAC and the
second generation of racecars adopted by the series, introduced in 2024 as the successor
of the AV-21. It is built by Dallara on the chassis of the Indy NXT, the open-wheel
single-seater of the development series of IndyCar, and has been adapted to driverless
operation while keeping the layout and the performance of the original car. The
vehicle has been designed to compete on oval tracks.

To enable autonomous operation, the cockpit hosts the on-board computer and the set of
actuators that turns the commands of the control software into steering, throttle,
brake and gear inputs to the vehicle, together with the power and network distribution
that connects sensors, computer and actuators. The bodywork has been modified to carry
the sensors, described in Section~\ref{sec:sens_av24}, and a radio link connects the
car to the base station, which receives the telemetry and sends the race control
signals. The main characteristics of the vehicle are collected in
Table~\ref{tab:av24}.

\begin{figure}[tb]
    \centering
    \figplaceholder{0.75\textwidth}{5cm}{Dallara AV-24 \TODO{photo}}
    \caption{The Dallara AV-24. \TODO{location and date}}
    \label{fig:av24}
\end{figure}

\begin{table}[tb]
    \centering
    \small
    \begin{tabular}{p{0.3\textwidth}p{0.45\textwidth}}
        \toprule
        Item              & Specification \\
        \midrule
        Series            & \\
        Base chassis      & \\
        Power unit        & \\
        Gearbox           & \\
        Top speed         & \\
        Mass              & \\
        Length / width / height & \\
        On-board computer & \\
        \bottomrule
    \end{tabular}
    \caption{Main characteristics of the Dallara AV-24. \TODO{to be completed}}
    \label{tab:av24}
\end{table}

% =====================================================================
\section{Dallara EAV-24}
\label{sec:eav24}

The Dallara EAV-24 (Fig.~\ref{fig:eav24}) is the official vehicle of the A2RL. It is
developed from the Dallara SF23, the single-seater of the Japanese Super Formula
championship, and keeps its chassis and aerodynamic package.

Unlike the AV-24, the EAV-24 is conceived for road circuits; therefore, the vehicle is
designed for high performance in heavy braking, in slow and fast corners and, in
general, in conditions of changing dynamics. The power unit is a turbocharged
\SI{2.0}{l} four-cylinder engine delivering about \SI{550}{hp}, coupled with a
longitudinal six-speed sequential gearbox, which brings the car to a top speed of about
\SI{250}{km/h}.

As on the AV-24, the cockpit hosts the on-board computer and the actuation system that
executes the commands of the control software, and the bodywork carries the sensors
described in Section~\ref{sec:sens_eav24}. The main characteristics are summarized in
Table~\ref{tab:eav24}.

\begin{figure}[tb]
    \centering
    \figplaceholder{0.75\textwidth}{5cm}{Dallara EAV-24 \TODO{photo}}
    \caption{The Dallara EAV-24. \TODO{location and date}}
    \label{fig:eav24}
\end{figure}

\begin{table}[tb]
    \centering
    \small
    \begin{tabular}{p{0.3\textwidth}p{0.45\textwidth}}
        \toprule
        Item              & Specification \\
        \midrule
        Series            & \\
        Base chassis      & \\
        Power unit        & \\
        Gearbox           & \\
        Top speed         & \\
        Mass              & \\
        Length / width / height & \\
        Wheelbase         & \\
        On-board computer & \\
        \bottomrule
    \end{tabular}
    \caption{Main characteristics of the Dallara EAV-24. \TODO{to be completed}}
    \label{tab:eav24}
\end{table}

% =====================================================================
\section{Sensors}
\label{sec:sensors}

This section describes the sensors mounted on the Dallara AV-24 and EAV-24. For each
of them, the working principle, the main advantages and the main limitations are
briefly described. Overall, the set of sensors allows a complete reconstruction of the
environment around the vehicle, within the distances of interest for racing, and
provides the reliable perception needed for autonomous driving.

\subsection{LiDAR}
\label{subsec:lidar}

A LiDAR (Light Detection and Ranging) is an active sensor that emits laser pulses and
measures the distance of the surfaces that reflect them from the time of flight of each
pulse, $r = c\,\Delta t/2$. A scanning mechanism steers the beam in azimuth and
elevation, and the high pulse repetition rate allows hundreds of thousands of points to
be acquired every second. The output of each scan is
a point cloud: a dense three-dimensional sampling of the surrounding surfaces, with
centimetre-level range precision and an angular resolution of a few hundredths of a
degree. From the point cloud the shape and the size of the objects around the vehicle
can be reconstructed; in the perception module it is processed to extract the bounding
boxes and the detections of the other vehicles (Section~\ref{sec:sw_arch}).

The long-range LiDARs mounted on both vehicles operate at a wavelength of
\SI{1550}{nm}, which allows a range of around \SI{80}{m}. Many units also allow the scan pattern to be configured, for instance by
concentrating the scan lines in a region of interest around the horizon, where the
other vehicles are expected.

Due to the nature of its working principle, the quality of the measurements degrades in
adverse weather, since rain, fog and the spray raised by other cars attenuate and
scatter the laser beam. Moreover, the energy returned by a target decreases quickly
with the distance and with the reflectivity of its surface, so that the working range
is shorter than that of RaDAR technologies. Finally, a LiDAR measures positions only:
the velocity of a target is not observed directly and has to be inferred from
successive scans.

\subsection{RaDAR}
\label{subsec:radar}

A RaDAR (Radio Detection and Ranging) is an active sensor that emits electromagnetic
waves in the radio band and receives the echoes reflected by the targets. Automotive
RaDARs are typically frequency-modulated continuous-wave (FMCW) sensors: the
transmitted signal is a sequence of frequency ramps, and the frequency difference
between the transmitted and the received signal is proportional to the range of the
target. The direction of the target is obtained from the phase differences of the echo
across the receiving antennas.

The main strength of the RaDAR is the direct measurement of velocity. When the target
moves with respect to the sensor, the phase of its echo changes from one ramp to the
next by an amount proportional to the radial component of the relative velocity
(Doppler effect):
\begin{equation}
    \Delta\varphi = \frac{4\pi}{\lambda}\,\dot\rho\,T_c
    \quad\Longrightarrow\quad
    \dot\rho = \frac{\lambda\,\Delta\varphi}{4\pi\,T_c},
    \label{eq:doppler}
\end{equation}
where $\lambda$ is the wavelength of the carrier, $T_c$ the time between two consecutive
ramps and $\Delta\varphi$ the phase difference of the echo between them; equivalently,
the echo is shifted in frequency by $|f_D| = 2|\dot\rho|/\lambda$ (Doppler shift). The
RaDAR therefore measures the \emph{range rate} $\dot\rho$, the rate of change of the
distance between the two vehicles, within a single scan. This is a unique property
among the sensors on board: LiDARs and cameras observe only positions, and the velocity
of a target can be obtained from them only by differentiating successive measurements.
The second strength is the range: radio waves are much less attenuated than light, so
the RaDAR detects a vehicle at larger distances than the LiDAR and extends the distance
at which an opponent is visible. For the same reason, it is also much less sensitive to
adverse weather, and its operation is guaranteed both by day and by night.

The quality of the position measurement, on the other hand, is lower than that of the
LiDAR. The RaDAR does not reconstruct the three-dimensional shape of the target: the
echoes come from a few points on the most reflective parts of the vehicle. In addition,
the angular resolution is considerably coarser than that of LiDARs and cameras.

\subsection{Cameras}
\label{subsec:camera}

Cameras are passive sensors that capture images of the scene at high resolution and
frame rate. The richness of the information contained in the images, such as colour
and texture, allows neural networks to detect and classify the surrounding objects
reliably: for each vehicle in the field of view, the detector returns a bounding box in
the image. The three-dimensional position of the vehicle is then reconstructed by
fitting its bounding box, combining the information coming from multiple cameras with
the knowledge of their calibration.

The main limitation of this approach is the possible degradation of the depth
reconstruction. Since the distance of the vehicle is extracted from the dimensions of
its bounding box, small errors in the size or in the position of the box in the image
translate into errors in the reconstructed depth, especially for distant targets.

\subsection{Localization}
\label{subsec:gnss}

The ego vehicle localizes itself with RTK-GNSS receivers coupled with inertial
measurement units, which provide position, attitude, velocity and acceleration in the
map frame with centimetre-level position accuracy. The ego state is used by the
perception pipelines to express the detections in the map frame and by the tracker to
compensate the motion of the ego vehicle.

\subsection{Dallara AV-24}
\label{sec:sens_av24}

The AV-24 is equipped with three Luminar Iris LiDARs, one facing forward and two on the
sides, which cover the front and the lateral regions of the car. Two Continental RaDARs
complete the perception of the surroundings: one on top of the cockpit, facing forward,
and one below the rear wing, facing backward, which covers the region behind the car.
Six Allied Vision Mako G319C cameras, arranged as a front stereo pair, two side cameras
and two cameras on the roll hoop, provide visual coverage around the vehicle.
Localization relies on a VectorNav VN-310 GNSS/INS and two NovAtel PwrPak7 receivers,
connected to a set of GNSS antennas distributed on the car. The sensors are listed in
Table~\ref{tab:sens_av24}, and their position on the car is shown in
Fig.~\ref{fig:av24_layout}.

\begin{table}[tb]
    \centering
    \small
    \begin{tabular}{>{\raggedright\arraybackslash}p{0.12\textwidth}>{\raggedright\arraybackslash}p{0.19\textwidth}>{\raggedright\arraybackslash}p{0.13\textwidth}>{\raggedright\arraybackslash}p{0.41\textwidth}}
        \toprule
        Sensor       & Model & Units & Main specifications \\
        \midrule
        LiDAR        &       &       & \\
        RaDAR        &       &       & \\
        Camera       &       &       & \\
        GNSS/INS     &       &       & \\
        \bottomrule
    \end{tabular}
    \caption{Sensor configuration of the Dallara AV-24. \TODO{to be completed}}
    \label{tab:sens_av24}
\end{table}

\begin{figure}[tb]
    \centering
    \figplaceholder{0.85\textwidth}{5.5cm}{Top view of the AV-24 with position and
    horizontal FoV of the sensors \TODO{draw}}
    \caption{Sensor layout of the Dallara AV-24.}
    \label{fig:av24_layout}
\end{figure}

\subsection{Dallara EAV-24}
\label{sec:sens_eav24}

The EAV-24 is equipped with three Seyond Falcon K LiDARs, one facing forward and two
on the sides of the car, and four ZF ProWave RaDARs, placed at the front, on the two
sides and at the rear, which together provide full \SI{360}{\degree} coverage. Seven
Leopard Imaging cameras complete the perception suite: two at the front, two on the
sides and three at the rear. Localization relies
on a dual-antenna VectorNav VN-310 RTK-GNSS/INS, which includes an inertial measurement
unit, and on a Kistler SF~Motion optical sensor, which measures the longitudinal and
lateral ground speed of the vehicle. Table~\ref{tab:sens_eav24} summarizes the
configuration and Fig.~\ref{fig:eav24_layout} shows the position of the sensors.

\begin{table}[tb]
    \centering
    \small
    \begin{tabular}{>{\raggedright\arraybackslash}p{0.12\textwidth}>{\raggedright\arraybackslash}p{0.19\textwidth}>{\raggedright\arraybackslash}p{0.13\textwidth}>{\raggedright\arraybackslash}p{0.41\textwidth}}
        \toprule
        Sensor       & Model & Units & Main specifications \\
        \midrule
        LiDAR        &       &       & \\
        RaDAR        &       &       & \\
        Camera       &       &       & \\
        GNSS/INS     &       &       & \\
        Ground speed &       &       & \\
        \bottomrule
    \end{tabular}
    \caption{Sensor configuration of the Dallara EAV-24. \TODO{to be completed}}
    \label{tab:sens_eav24}
\end{table}

\begin{figure}[tb]
    \centering
    \figplaceholder{0.85\textwidth}{5.5cm}{Top view of the EAV-24 with position and
    horizontal FoV of the sensors \TODO{draw}}
    \caption{Sensor layout of the Dallara EAV-24.}
    \label{fig:eav24_layout}
\end{figure}

% =====================================================================
\section{Software Architecture}
\label{sec:sw_arch}

The autonomous driving software, shared by the two vehicles, is organized in four
main modules: localization, perception, planning and control (Fig.~\ref{fig:sw_arch}).
Each module runs continuously and exchanges data with the others, so that the
information flows from the sensors to the actuators through a chain in which the
output of one module is the input of the next.

\begin{figure}[tb]
    \centering
    \figplaceholder{0.95\textwidth}{6cm}{Block diagram: sensors $\rightarrow$
    perception (detection $\rightarrow$ target tracking) $\rightarrow$ planning
    $\rightarrow$ control $\rightarrow$ actuators; localization feeding perception,
    planning and control; track database feeding planning \TODO{draw}}
    \caption{Main modules of the autonomous driving software and data exchanged between
    them.}
    \label{fig:sw_arch}
\end{figure}

\paragraph{Localization.} The localization module estimates the state of the ego
vehicle, i.e.\ its position, heading, velocity and acceleration in the map frame, by
fusing the measurements of the GNSS receivers, of the inertial units and of the
proprioceptive sensors. Its output is used by the perception, planning and control
modules.

\paragraph{Perception.} The perception module turns the raw data of LiDARs, RaDARs and
cameras into a description of the vehicles around the ego car. A first stage processes
the data of each sensor separately: point clouds, RaDAR reflections and images are
filtered, the points belonging to other vehicles are isolated and a list of
\emph{detections} is produced, each one containing the position of a vehicle and, for
the RaDAR, its range rate. The detections of the different sensors are then passed to
the \emph{target tracking} module, which associates them over time with the vehicles
already known, fuses the information coming from the different sensors and estimates
the state of each opponent, including quantities that are not measured directly. The result is
the list of the opponents with their estimated state, which is the output of perception towards planning. The target
tracking module is the subject of this thesis and is described in detail in
Chapter~\ref{ch:tracking}.

\paragraph{Planning.} It receives the ego state from localization, the opponent list
from perception and the information stored in a track database, which contains the
boundaries of the track and a set of reference trajectories. Based on the presence of
other vehicles and on the chosen strategy, i.e.\ following, overtaking or defending,
the planning module decides the trajectory that the vehicle has to follow; the
trajectory can also be selected by configuration. The output is a reference trajectory,
i.e.\ a path with the associated speed profile, which is sent to control.

\paragraph{Control.} It receives the reference trajectory from planning and the ego
state from localization, and computes the commands for the actuators through two
controllers: a longitudinal controller, which acts on throttle, brake and gear to track
the reference speed, and a lateral controller, which computes the steering angle that
keeps the vehicle on the reference path.

% =====================================================================
\section{Datasets}
\label{sec:datasets}

% =====================================================================
\section{Ground Truth Reconstruction}
\label{sec:gt}
